\subsubsection{Ferramentas de software e bibliotecas}

O Buildroot foi selecionado como toolchain de cross-compilação para o desenvolvimento de um sistema operacional personalizável baseado em Linux, devido ao desempenho demonstrado em trabalhos correlatos, como o de \citeonline{ozturk2023}. Para viabilizar o projeto de uma Placa de Circuito Impresso (PCB, do inglês \textit{Printed Circuit Board}), foi selecionado o Autodesk Fusion, uma vez que possui a ferramenta de Automação de Projetos Eletrônicos (EDA, do inglês \textit{Electronic Design Automation}), além de possuir uma curva de aprendizado rápida, interface intuitiva, licenciamento gratuito e um desenvolvimento colaborativo em nuvem.

Para o desenvolvimento dos algoritmos de ML, foi adotada a linguagem de programação Python por ser amplamente consolidada nesse domínio por dispor de um ecossistema robusto de bibliotecas especializadas, que oferecem implementações otimizadas de modelos, pipelines de treinamento e ferramentas de avaliação. Ademais, a sintaxe simples e expressiva do Python favorece a prototipagem rápida e a legibilidade do código, aspectos relevantes em projetos de pesquisa que demandam iterações frequentes sobre os modelos desenvolvidos.

Por fim, para o desenvolvimento da interface gráfica embarcada executada diretamente no hardware, foi selecionado o \textit{framework} Qt em conjunto com a linguagem C++. O Qt é uma plataforma madura e amplamente utilizada no desenvolvimento de interfaces gráficas para sistemas embarcados, oferecendo suporte a renderização eficiente, portabilidade entre diferentes arquiteturas e integração nativa com o sistema operacional Linux, adotado neste trabalho \cite{qtAbout}. A escolha do C++ como linguagem base justifica-se pelo seu desempenho superior em ambientes com recursos computacionais limitados, onde o controle fino sobre memória e tempo de execução é determinante para a responsividade da interface \cite{knuttila2012functional}.


% O Buildroot foi selecionado como toolchain de cross-compilação para o desenvolvimento de um sistema operacional personalizável baseado em Linux, devido ao desempenho demonstrado em trabalhos correlatos, como o de \citeonline{ozturk2023}. Trata-se de um framework amplamente adotado na comunidade embarcada por permitir a geração de sistemas de arquivos raiz, kernels e toolchains de forma automatizada e reprodutível, sendo especialmente adequado para plataformas com recursos computacionais limitados.

% Adicionalmente, para viabilizar o projeto de uma Placa de Circuito Impresso (PCB, do inglês \textit{Printed Circuit Board}), foi selecionado um software dotado de ferramenta de Automação de Projetos Eletrônicos (EDA, do inglês \textit{Electronic Design Automation}). Esse tipo de ferramenta permite conduzir o ciclo completo de projeto eletrônico, abrangendo a captura de esquemáticos, simulação de circuitos, verificação de regras de projeto (DRC, do inglês \textit{Design Rule Check}) e definição do layout físico da placa \cite{autodeskEda}.

% A escolha de uma ferramenta EDA adequada é determinante para a confiabilidade do hardware desenvolvido, uma vez que erros identificados ainda na fase de simulação ou verificação reduzem significativamente o custo e o tempo de correção em etapas posteriores de fabricação. A Tabela~\ref{tableSoftwareEBAComparer} apresenta um comparativo entre as principais ferramentas de software para desenvolvimento de PCB consideradas neste trabalho, avaliando critérios como custo de licenciamento, curva de aprendizado, qualidade dos recursos de simulação e suporte à colaboração em nuvem.
    
% \begin{table}[h]
% \caption{Comparativo entre Softwares com ferramentas EDA do mercado}
% \label{tableSoftwareEBAComparer}
% \centering
% \begin{tabular}{ | m{8em} | m{7em} | m{7em} | m{5em} | m{6em} | }
%     \hline
%     \textbf{Software} & \textbf{Preço} & \textbf{Curva de aprendizado} & \textbf{Simulação} & \textbf{Colaboração  em nuvem} \\
%     \hline
%     Proteus                     & \$250,00 - \$3000      & Média             & Forte     & Não \\
%     OrCAD / Cadence             & Variável por demanda  & Alta              & Forte     & Sim \\
%     \textbf{Autodesk Fusion}    & Gratuito              & Rápida - Média    & Básica    & Sim \\
%     \hline
% \end{tabular}
% \footnotesize \textbf{Fonte: Elaborado pelos autores.}
% \end{table}

% Apesar de o software Autodesk Fusion não ser referência em simulação de PCB, ele se destaca entre as demais opções analisadas por reunir um conjunto de características adequadas ao contexto deste trabalho. A ferramenta apresenta uma curva de aprendizado reduzida, interface intuitiva e licenciamento gratuito, o que a torna acessível sem impacto no orçamento do projeto. Além disso, disponibiliza recursos nativos de desenvolvimento colaborativo em nuvem, um diferencial particularmente relevante neste contexto, uma vez que o projeto é desenvolvido por mais de um autor. Dessa forma, as limitações da ferramenta no que se refere à simulação não representam um fator crítico, considerando a baixa complexidade do hardware projetado. Portanto, os recursos de verificação e DRC oferecidos pelo Autodesk Fusion são suficientes para garantir a confiabilidade do layout produzido.

% Para o desenvolvimento dos algoritmos de ML, foi adotada a linguagem de programação Python, amplamente consolidada nesse domínio por dispor de um ecossistema robusto de bibliotecas especializadas, como Scikit-learn, TensorFlow e PyTorch, que oferecem implementações otimizadas de modelos, pipelines de treinamento e ferramentas de avaliação. Ademais, a sintaxe simples e expressiva do Python favorece a prototipagem rápida e a legibilidade do código, aspectos relevantes em projetos de pesquisa que demandam iterações frequentes sobre os modelos desenvolvidos.

% Para o desenvolvimento da interface gráfica embarcada executada diretamente no hardware, foi selecionado o \textit{framework} Qt em conjunto com a linguagem C++. O Qt é uma plataforma madura e amplamente utilizada no desenvolvimento de interfaces gráficas para sistemas embarcados, oferecendo suporte a renderização eficiente, portabilidade entre diferentes arquiteturas e integração nativa com o sistema operacional Linux, adotado neste trabalho \cite{qtAbout}. A escolha do C++ como linguagem base justifica-se pelo seu desempenho superior em ambientes com recursos computacionais limitados, onde o controle fino sobre memória e tempo de execução é determinante para a responsividade da interface \cite{knuttila2012functional}.


